%%%%%%%%%%%%%%%%%%%%%%%%%%%%%%%%%%%%%%%%%%%%%%%%%%%%%%%%%%%%%%%%%%%%%%
% How to use writeLaTeX: 
%
% You edit the source code here on the left, and the preview on the
% right shows you the result within a few seconds.
%
% Bookmark this page and share the URL with your co-authors. They can
% edit at the same time!
%
% You can upload figures, bibliographies, custom classes and
% styles using the files menu.
%
%%%%%%%%%%%%%%%%%%%%%%%%%%%%%%%%%%%%%%%%%%%%%%%%%%%%%%%%%%%%%%%%%%%%%%

\documentclass[12pt]{article}

\usepackage{sbc-template}

\usepackage{graphicx,url}

%\usepackage[brazil]{babel}   
\usepackage[utf8]{inputenc}  

     
\sloppy

\title{Exercício 3 Sistemas Distribuídos}

\author{Christian Aguiar Plentz\inst{1}, João Roberto Lemos Fidellis\inst{2} }


\address{Ciência da Computação -- Universidade do Vale do Rio dos Sinos (UNISINOS)\\
  Caixa Postal 275 -- 93022-750 -- São Leopoldo -- RS -- Brazil
  \email{\{plentzchristian,JRFIDELLIS\}@edu.unisinos.br}
}


\begin{document} 

\maketitle

\section{ Escreva um programa em Sockets Ping-Pong para medir o tempo
    de comunicação entre 2 processos interligados por uma rede LAN.
    não pode ser VM ou na mesma maquina}

Use os protocolos UDP e TCP e teste com os seguintes tamanhos de mensagem
\textbf{100B, 500B, 1KB, 100KB, 500KB, 1MB e 10MB}

\subsection{UDP}
    Os processos UDP conseguiram se comunicar ate os 1KB pois quando a mensagem
    é maior que 65507 Bytes (Tamanho máximo de um payload UDP) o próprio SO retorna
    uma mensagem de erro com o numero 90 nomeado \textit{EMSGSIZE} \\
\textbf{Saída do terminal:} \\
\indent \texttt{OSError: [Errno 90] Message too long}

\noindent \textbf{Tempo da comunicação entre os 2 processos}

\begin{center}
\begin{tabular}{|l|l|}  % | for vertical lines
\hline% Horizontal line
Tamanho do Pacote & Tempo em Milisegundos \\
\hline
    100B & 1.66ms \\
    500B & 1.40ms \\
    1KB & 1.35ms \\
\hline
\end{tabular}
\end{center}

\subsection{TCP}

Os processos TCP conseguiram se comunicar ate os 1MB, ao tentar os 10MB sempre houve problema de
timeout entre os pacotes. \\
\textbf{Saída do Terminal:} \\
\indent \texttt{TimeoutError: timed out}

\noindent \textbf{Tempo da comunicação entre os 2 processos:} \\
\begin{center}
\begin{tabular}{|l|l|}  % | for vertical lines
\hline% Horizontal line
Tamanho do Pacote & Tempo em Milisegundos \\
\hline
100B & 1.77ms \\
500B & 1.52ms \\
1KB & 1.30ms \\
100KB & 1.69ms \\
500KB & 1.97ms \\
1MB & 1.57ms \\
\hline
\end{tabular}
\end{center}

\clearpage

\section{Usando o mesmo programa anterior, faca o gráfico de uso de largura de banda.
    Explique o comportamento}


\subsection{Gráfico UDP}

\begin{figure}[ht]
\centering
\includegraphics[width=1\textwidth]{exercicio3/grafico_udp.png}
\end{figure}

Devido a somente ter 3 resultados pelas limitações do tamanho de uma mensagem UDP
não se consegue tirar muitas conclusões do gráfico além de que o UDP apresenta
uma largura de banda linear ao tamanho do pacote contanto que ele esteja nos limites UDP

\clearpage

\subsection{Gráfico TCP}

\begin{figure}[ht]
\centering
\includegraphics[width=1\textwidth]{exercicio3/grafico_tcp.png}
\end{figure}

O grafico foi usado uma sequencia logaritimica para medir seu tamanho devido aos números muito
altos que se teve no tamanho do pacote. mas mesmo o gráfico estando em sequencia logarítmica
a largura de banda também apresenta uma progressão linear

\end{document}

